% TOP_PROBLEM: {"id": 1277, "title": "Montgomery's Pair Correlation Conjecture", "category_id": 1, "category": {"id": 1, "name": "number_theory", "display_name": "Number Theory", "description": "Properties of integers, prime numbers, Diophantine equations.", "slug": "number-theory", "order_index": 1, "created_at": "2026-07-31T15:26:25.670Z"}, "status": "open"}
% FIRST_POSTED: 2026-09-27
% !TeX program = lualatex
% Compile twice: lualatex 1277.tex
% Standalone research notebook; original section preserved.
% Research additions: 27 September 2026. No external TeX input or BibTeX file.
\documentclass[11pt,a4paper]{article}
\usepackage[margin=24mm,headheight=14pt]{geometry}
\usepackage{fontspec}
\setmainfont{Latin Modern Roman}
\setsansfont{Latin Modern Sans}
\setmonofont{Latin Modern Mono}
\usepackage{amsmath,amssymb,mathtools}
\usepackage{unicode-math}
\setmathfont{Latin Modern Math}
\usepackage{microtype}
\usepackage{enumitem,xurl,needspace}
\usepackage[dvipsnames]{xcolor}
\usepackage{hyperref}
\hypersetup{unicode=true,colorlinks=true,linkcolor=MidnightBlue,urlcolor=MidnightBlue,
 pdftitle={Research notebook - Problem 1277},pdfauthor={Open Problems Math},bookmarksnumbered=true}
\usepackage{bookmark}
\usepackage{titlesec}
\titleformat{\section}{\Large\bfseries\raggedright}{\thesection}{0.7em}{}
\usepackage{fancyhdr}
\pagestyle{fancy}\fancyhf{}
\fancyhead[L]{\small\sffamily Open Problems Math | Research notebook}
\fancyhead[R]{\small\sffamily Problem 1277 | 27 September 2026}
\fancyfoot[C]{\thepage}
\setlength{\parindent}{0pt}
\setlength{\parskip}{5pt plus 1pt}
\setlength{\emergencystretch}{3em}
\setcounter{tocdepth}{1}
\setcounter{secnumdepth}{1}
\setlist[itemize]{leftmargin=3em,itemsep=5pt,topsep=4pt,parsep=0pt}
\allowdisplaybreaks
\numberwithin{equation}{section}
\newcommand{\N}{\mathbb{N}}
\newcommand{\Z}{\mathbb{Z}}
\newcommand{\Q}{\mathbb{Q}}
\newcommand{\R}{\mathbb{R}}
\newcommand{\C}{\mathbb{C}}
\newcommand{\F}{\mathbb{F}}
\newcommand{\A}{\mathbb{A}}
\newcommand{\PP}{\mathbb P}
\newcommand{\E}{\mathbb{E}}
\newcommand{\eps}{\varepsilon}
\newcommand{\dd}{\,\mathrm d}
\newcommand{\sref}[2]{\hyperlink{source:#1:#2}{[S#2]}}
\newcommand{\eref}[2]{\hyperlink{extra:#1:#2}{[E#2]}}
\newif\ifshowcatalogue
\showcataloguetrue
\newcommand{\cataloguescope}[1]{\ifshowcatalogue
 \paragraph{Exact scope in the supplied catalogue.}
 \begin{quote}\small #1\end{quote}\fi}
\begin{document}
\setcounter{section}{0}
\section*{\texorpdfstring{Montgomery's Pair Correlation Conjecture}{Montgomery's Pair Correlation Conjecture}}\label{1277}
\subsection{Definitions and mathematical statement}
List the positive ordinates $\gamma$ of nontrivial zeta zeros with their multiplicities, and write $N(T)=\#\{0<\gamma\le T\}$. A standard fixed-separation form of the conjecture is, for every $0<a<b$,
\[
 \frac{1}{N(T)}\#\left\{0<\gamma,\gamma'\le T:
 a\le\frac{\log T}{2\pi}(\gamma-\gamma')\le b\right\}
 \longrightarrow\int_a^b\left(1-\left(\frac{\sin\pi u}{\pi u}\right)^2\right)du.
\]
Only positive separation is counted, so diagonal terms are absent. The scale $2\pi/\log T$ is the mean vertical spacing near height $T$. This is a pair-statistics law, not a claim that consecutive spacings have this density.
\par\smallskip\textit{Formulation and concept references:} \sref{1277}{1}.
\cataloguescope{Prove that normalized vertical gaps between pairs of nontrivial zeros of the Riemann zeta function have limiting pair density 1 - (sin(pi u)/(pi u))\textasciicircum{}2 at fixed positive normalized separation.}
\subsection{Short English statement}
Do pairs of zeta zeros have the universal repulsion statistics predicted by the sine-kernel law?
% BEGIN RESEARCH_ADDITION ATTEMPT_74
\subsection{Research attempt}

\paragraph{Current frontier and quantifiers.}
The source's Section~4 was inspected: the target concerns ordinates and does
not assume RH. Goldston--Lee--Schettler--Suriajaya show that their PCC
implies density-one simplicity and density-one location on the critical
line, rather than RH for every zero \eref{1277}{1}.
Restricted Fourier-support theorems are the foundational starting point,
but the target tests every fixed positive interval. The attempt below
isolates the analytic information lost by restricting Fourier support.
It does not conflate pair statistics with consecutive gaps.

\paragraph{Attack route.}
An explicit-formula route needs new control of off-diagonal prime sums.
A positivity route hopes that the known low-frequency information fixes
the whole pair law. A measure-convergence route asks for a uniform tail
estimate that upgrades Fourier tests to interval counts. The second route
is actively falsified below; the third is then stated as a precise
sufficient package. No unproved prime correlation estimate is smuggled
into that package as an established bound.

\paragraph{Attempt: the exact Fourier measure.}
Let $c_T=\log T/(2\pi)$ and include all ordered pairs to form
\[
 \nu_T=\frac1{N(T)}\sum_{0<\gamma,\gamma'\leq T}
                          \delta_{c_T(\gamma-\gamma')}.
\]
At zero this measure contains diagonal and repeated-ordinate contributions.
For a Schwartz test function $h$, with
$\widehat h(\xi)=\int h(u)e^{-2\pi i\xi u}\,du$, finite summation and
Fourier inversion give
\begin{equation}\label{eq74:fourier}
 \int h\,d\nu_T=\int_{\R}\widehat h(\xi)F_T(\xi)\,d\xi,
 \qquad
 F_T(\xi)=\frac1{N(T)}
       \left|\sum_{0<\gamma\leq T}e^{2\pi i\xi c_T\gamma}\right|^2.
\end{equation}
The sign of $\xi$ is immaterial because $F_T$ is even. In particular
$F_T\geq0$, but its trivial pointwise bound $N(T)$ grows with $T$.
The target itself asks only for convergence of $\nu_T([a,b])$ when
$0<a<b$; it does not by definition specify the zero atom.

A sufficient, possibly stronger, full-measure candidate is
\[
 \nu_\infty=\delta_0+g(u)\,du,
 \qquad g(u)=1-\left(\frac{\sin\pi u}{\pi u}\right)^2.
\]
Its Fourier transform is the positive tempered measure
$\delta_0+\min(|\xi|,1)\,d\xi$, since the Fourier transform of the squared
sinc is $(1-|\xi|)_+$. Stating this sufficient candidate does not assume
that the original interval formulation alone has already supplied all
necessary tightness and zero-mass information.

\paragraph{Attempt: low-frequency information does not determine $g$.}
Fix any finite frequency cutoff $A>0$. Choose a nonzero even real smooth
function $b$ of compact support entirely in $|\xi|>A+1$. Let $\psi$ be its
inverse Fourier transform and set $h(u)=u^2\psi(u)$. Then $h$ is real,
even and Schwartz, and $\widehat h=-b''/(4\pi^2)$ also vanishes on
$[-A,A]$. Near zero,
\[
 g(u)=\frac{\pi^2}{3}u^2+O(u^4),\qquad h(u)=O(u^2).
\]
Away from zero, $g$ is strictly positive; at infinity it tends to one
while $h$ decays. Therefore $|h(u)|\leq K g(u)$ for some finite $K$.
For $0<\varepsilon<1/K$, both
\begin{equation}\label{eq74:perturb}
 \delta_0+(g+\varepsilon h)\,du,
 \qquad \delta_0+(g-\varepsilon h)\,du
\end{equation}
are positive measures. They agree on every Schwartz test whose Fourier
transform is supported in $[-A,A]$, by Fourier inversion and the support
condition on $\widehat h$. They differ on some interval in $(0,\infty)$:
$h$ is not identically zero there, so continuity gives an interval with
nonzero integral.

Even positivity of the Fourier transform can be retained. On the support
of $\widehat h$ the candidate Fourier density is one, and $\widehat h$
is bounded. Reducing $\varepsilon$ so that
$\varepsilon\|\widehat h\|_\infty<1$ makes
$\delta_0+[\min(|\xi|,1)\pm\varepsilon\widehat h(\xi)]d\xi$
positive as well. Thus positivity on both sides and any one fixed Fourier
window cannot by themselves determine the sine-kernel pair measure.
This is an analytic countertest to a uniqueness argument. The perturbed
measures are not claimed to be pair measures of point processes, much
less of zeta zeros.

\paragraph{Attempt: a sufficient tail-control upgrade.}
Suppose the following two additional assertions held for the actual
$F_T$. First, for every smooth compactly supported $\varphi$,
\begin{equation}\label{eq74:local}
 \int\varphi(\xi)F_T(\xi)\,d\xi\longrightarrow
 \varphi(0)+\int\varphi(\xi)\min(|\xi|,1)\,d\xi.
\end{equation}
Second, for fixed constants $B,C$ independent of $T$ and every integer $j$,
\begin{equation}\label{eq74:tail}
 \int_j^{j+1}F_T(\xi)\,d\xi\leq C(1+|j|)^B.
\end{equation}
Then \eqref{eq74:local} extends from compactly supported functions to
Schwartz functions. To prove this, multiply a Schwartz function by a
smooth cutoff equal to one on $[-R,R]$. Its omitted integral is bounded,
uniformly in $T$, by
$C\sum_{|j|\geq R-1}(1+|j|)^B\sup_{[j,j+1]}|\varphi|$.
Rapid decay makes this tail tend to zero. The limiting measure has the
same property, so taking $T\to\infty$ first and then $R\to\infty$ is
justified. Equation \eqref{eq74:fourier} now proves convergence on every
smooth compactly supported $h$ in physical space.

For $0<a<b$, choose smooth minorants and majorants for
$\mathbf1_{[a,b]}$ supported away from zero, whose differences have
arbitrarily small integral against the continuous density $g$. Positivity
of $\nu_T$ sandwiches the liminf and limsup. Sending the approximation
error to zero gives the required interval limit. Thus the two displayed
assertions suffice for the exact stated conclusion.

\paragraph{Stress test.}
The missing part of \eqref{eq74:local} is precisely the unrestricted
frequency range. A theorem for one fixed window is not the quantified
assertion for every compact support. The perturbations
\eqref{eq74:perturb} show why no positivity shortcut removes this gap.
Likewise $F_T\leq N(T)$ does not supply \eqref{eq74:tail} with constants
independent of height, so dominated convergence cannot be invoked with
that bound. The sharp cutoff in $\nu_T$, any smoothed zero weights used in
an explicit formula, and the absence of an RH assumption must all be
matched before importing a Fourier theorem.

The full-measure sufficient package includes a unit zero atom, stronger
information than merely declaring that diagonal pairs are excluded from
$[a,b]$. It has not been extracted for free from the catalogue's statement.
Nor does the density $g$ describe consecutive spacings: each zero is paired
with every other zero. The argument neither computes zeros nor infers a
limiting law from a finite plot.

\paragraph{Outcome.}
\textbf{Outcome: Exploratory approach with unresolved gap.}

An explicit positive, positive-Fourier perturbation disproves the proposed
uniqueness inference from low-frequency data. A proved measure-theoretic
upgrade identifies sufficient full-frequency convergence and uniform tail
control for the actual interval target. Those analytic inputs are not
established. No pair-correlation theorem, zero-location theorem, or
counterexample for zeta itself is obtained; the new reasoning is a
methodological test without a priority claim.

% Keep the local bibliography together rather than leave a source fragment.
\needspace{170mm}
% END RESEARCH_ADDITION ATTEMPT_74
\subsection{Sources}
\begin{itemize}\raggedright
\item[\textnormal{[S1]}]\hypertarget{source:1277:1}{} Goldston, Lee, Schettler, and Suriajaya, Pair Correlation Conjecture for the zeros of the Riemann zeta-function I (version 4, 2026), Section 4, PCC.
\url{https://arxiv.org/html/2503.15449v4}.
{\small\textit{Catalogue formulation source.}}
\item[\textnormal{[S2]}]\hypertarget{source:1277:2}{} Variations of the Hardy Z-Function and the Montgomery Pair Correlation Conjecture.
\url{https://arxiv.org/abs/2511.18275}.
{\small\textit{Catalogue research/status reference; not a proof endorsement.}}
\item[\textnormal{[S3]}]\hypertarget{source:1277:3}{} Pair correlation of zeros of Dirichlet L-functions: a possible path towards the conjectures of Chowla, Elliott-Halberstam and Montgomery.
\url{https://link.springer.com/article/10.1007/s00208-026-03383-y}.
{\small\textit{Catalogue research/status reference; not a proof endorsement.}}
% BEGIN RESEARCH_ADDITION REFERENCES_74
\item[\textnormal{[E1]}]\hypertarget{extra:1277:1}{} Daniel A. Goldston,
Junghun Lee, Jordan Schettler and Ade Irma Suriajaya,
\emph{Pair Correlation Conjecture for the zeros of the Riemann zeta-function I:
simple and critical zeros}, arXiv:2503.15449v4, Section 4 and Theorem 1.
\url{https://arxiv.org/html/2503.15449v4}.
{\small\textit{Inspected 27 September 2026. The source's PCC includes
intervals approaching zero through its stated counting formulation;
the present attempt keeps the supplied positive-interval target and uses
a separately labelled sufficient full-measure package.}}
% END RESEARCH_ADDITION REFERENCES_74
\end{itemize}

\end{document}
